\documentclass[12pt]{article}
\newcommand{\DocShort}{Ajuste lineal con funciones de Excel}
\input{preamble_population.tex}
\begin{document}
\doccover{Ajuste lineal por mínimos cuadrados con funciones de Excel}{2026-05}
\handoutintro{Ajustar una recta a los datos de población para proyectar con mínimos cuadrados.}{Series con tendencia aproximadamente lineal en el horizonte analizado.}

\section{Datos ejemplo}
Aquí se usa $x$ como tiempo medido en años desde 1970.

\begin{center}
\begin{tabular}{lcccccc}
\toprule
Año & 1970 & 1980 & 1990 & 2000 & 2010 & 2020 \\
\midrule
$x$ (años) & 0 & 10 & 20 & 30 & 40 & 50 \\
Población & 24,000 & 31,200 & 38,800 & 46,300 & 54,500 & 62,800 \\
\bottomrule
\end{tabular}
\end{center}

\section{Ecuaciones mínimas}
El modelo lineal es:
\begin{equation}
P=a+bx
\end{equation}
donde $P$ es la población proyectada, $a$ es la ordenada al origen, $b$ es la pendiente y $x$ es el tiempo en años desde el origen elegido.

En Excel:
\[
\texttt{=LINEST(rango\_P, rango\_x, TRUE, TRUE)}
\]
En Excel en español:
\[
\texttt{=ESTIMACION.LINEAL(rango\_P; rango\_x; VERDADERO; VERDADERO)}
\]

\excelnote{Organice $x$ en una columna y $P$ en otra. Use \texttt{LINEST} o \texttt{ESTIMACION.LINEAL} para obtener pendiente e intercepto. Luego evalúe la recta en el año deseado.

\textbf{Nota.} Las funciones \texttt{TREND} y \texttt{TENDENCIA} son útiles para proyectar valores, pero no devuelven directamente los parámetros de la ecuación.}

\section{Procedimiento}
\begin{enumerate}
\item Definir una variable temporal $x$.
\item Cargar $x$ y $P$ en Excel.
\item Aplicar \texttt{LINEST}.
\item Escribir la ecuación ajustada.
\item Sustituir el valor de $x$ para los años futuros.
\end{enumerate}

\section{Ejemplo resuelto}
Con los datos anteriores, \texttt{LINEST} (o \texttt{ESTIMACION.LINEAL}) devuelve aproximadamente:
\[
a=\num{23547.62},\qquad b=\num{775.43}\ \text{hab/año}
\]

Por tanto,
\[
P = \num{23547.62} + \num{775.43}x
\]

Para 2030, 2040 y 2050:
\[
P_{2030}=\num{23547.62}+\num{775.43}(60)=\num{70073}
\]
\[
P_{2040}=\num{23547.62}+\num{775.43}(70)=\num{77828}
\]
\[
P_{2050}=\num{23547.62}+\num{775.43}(80)=\num{85582}
\]

\resultbox{El ajuste lineal con LINEST produce 70,073 habitantes en 2030, 77,828 en 2040 y 85,582 en 2050.}

\section{Teoría breve}
El método de mínimos cuadrados minimiza la suma de los cuadrados de los errores verticales entre los datos observados y la recta ajustada.

\section{Observaciones de uso}
\begin{itemize}
\item Útil como alternativa objetiva al método aritmético.
\item Si la nube de puntos es curvada, un ajuste lineal puede sesgar la proyección.
\item Debe aclararse siempre la unidad temporal de $x$; en este ejemplo es años.
\end{itemize}
\end{document}
